\section{其他重正化方案}
\label{sec:ren}

尽管我们在推导上述匹配公式时假设准 PDF 在 $\overline{\ensuremath{\operatorname{MS}}}$ 方案下重正化，这并不是对我们结果的限制。
既然规范不变的 Wilson 线算符 $\tilde{O}_\Gamma(z)$ 已被证明在坐标空间可乘法重正化~\cite{Ji:2017oey,Ishikawa:2017faj}，那么在使用上述因子化公式之前，总可以把 $\tilde{Q}_{\Gamma}(z)$ 从任何其他方案转换到 $\overline{\ensuremath{\operatorname{MS}}}$ 方案。准 PDF 的重正化近年在许多文章中被研究过~\cite{Ji:2015jwa,Ishikawa:2016znu,Chen:2016fxx,Xiong:2017jtn,Constantinou:2017sej,Alexandrou:2017huk,Chen:2017mzz,Green:2017xeu,Stewart:2017tvs,Wang:2017eel}。我们将讨论其中的一些结果，并说明如何把它们纳入式~(\ref{eq:fac-tq})的因子化公式。

$\overline{\ensuremath{\operatorname{MS}}}$ 方案便于我们讨论 OPE，因为它保证洛伦兹与规范不变性，但它对格点重正化并不实用。由于格点理论有天然的紫外截断 $1/a$（$a$ 为格距），未重正化的\spcorr $\tilde{Q}$ 按照式~(\ref{eq:ren})从 Wilson 线自能继承幂次发散。对任意方案 $X$，重正化\spcorr
\begin{align}
\tilde{Q}^X(\zeta,z^2\mu^2_R)
  = \lim_{a\to0} Z^{-1}_X(z^2\mu_R^2,a^2\mu^2_R)\, \tilde{Q}(\zeta,z^2/a^2)
\end{align}
应当不含任何紫外发散并具有良好的连续极限 $a\to0$。特别地，这一连续极限与紫外调节无关，所以
\begin{align}
&\lim_{a\to0} Z^{-1}_X(z^2\mu_R^2,a^2\mu^2_R)\tilde{Q}(\zeta,z^2/a^2) \nn\\
&= Z^{-1}_X(z^2\mu^2_R,\epsilon)\tilde{Q}(\zeta,z^2,\epsilon)\,.
\end{align}
于是我们可以通过如下转换把 $\tilde{Q}^X(\zeta,z^2\mu^2_R)$ 与 $\overline{\rm MS}$ 方案联系起来：
\begin{align}
\tilde{Q}^X(\zeta,z^2\mu^2_R) =& \frac{Z_{\overline{\rm MS}}(\epsilon,\mu)}{Z_X(z^2\mu^2_R,\epsilon)}\, \tilde{Q}^{\overline{\rm MS}}(\zeta,z^2\mu^2)\nn\\
=& Z'_X(z^2\mu_R^2,\mu_R^2/\mu^2)\, \tilde{Q}^{\overline{\rm MS}}(\zeta,z^2\mu^2)\,,
\end{align}
其中调节器 $\epsilon$ 依赖在 $Z_{\overline{\rm MS}}$ 与 $Z_X$ 之间完全抵消。比值 $Z_X'$ 可以在 QCD 中微扰计算，文献~\cite{Constantinou:2017sej}对若干格点方案和 RI/MOM 方案做了这件事。于是我们在第~\ref{sec:ope}~节证明的因子化公式经系数函数的轻微修改后仍适用于 $\tilde{Q}^X$：
\begin{align}
\tilde{Q}^X(\zeta,z^2\mu^2_R)
 = \int_{-1}^1 \!\! d\alpha \
   \mathcal{C}^X(\alpha,\mu_R^2/\mu^2,\mu^2 z^2)\, Q(\alpha\zeta,\mu)\,,
\end{align}
其中方案 $X$ 的匹配系数与 $\overline{\rm MS}$ 的关系为
\begin{align}
\mathcal{C}^X(\alpha,\mu_R^2/\mu^2,\mu^2 z^2) = Z'_X(z^2\mu_R^2,\mu_R^2/\mu^2)\, \mathcal{C}(\alpha,\mu^2 z^2) \,.
\end{align}
对赝 PDF，修改后的结果也涉及这同一系数：
\begin{align}
\mathcal{P}^X(x,z^2\mu^2_R)
=&\int_{|x|}^1 {dy\over |y|}\
 \mathcal{C}^X\Bigl({x\over y},{\mu_R^2\over \mu^2},\mu^2 z^2\Bigr) q(y,\mu)
  \\
&+\int^{-|x|}_{-1} {dy\over |y|}\
 \mathcal{C}^X\Bigl({x\over y},{\mu_R^2\over \mu^2},\mu^2 z^2\Bigr) q(y,\mu)
 \,. \nn
\end{align}
而对准 PDF，我们有
\begin{align}\label{eq:rimomfac}
\tilde{q}_X\left(x,{\mu_R^2\over P_z^2}\right)
&\equiv \int \frac{d \zeta}{2 \pi}\: e^{ix \zeta}\:  \tilde{Q}^X \biggl(\zeta, \frac{{\mu}^2\zeta^2}{P_z^2}\biggr)
\\
&=\int_{-1}^1 {dy\over |y|}
\: C^X\Bigl(\frac{x}{y},\frac{\mu_R}{\mu},\frac{\mu}{|y| P^z}\Bigr)\, q(y,\mu)\,.\nn
\end{align}
这里 $X$ 方案的修改后系数与 $\overline{\rm MS}$ 方案系数的关系为
%
\begin{align}
 & C^X\Bigl(\frac{x}{y},  \frac{\mu_R}{\mu},\frac{\mu}{|y| P^z}\Bigr)
 \\
&\ =\!\! \int\!\! d\eta\ \bar{Z}'_X\Bigl(\eta^2, {\mu_R^2\over \mu^2}\Bigr)\:
 C\Bigl( \frac{x}{y} - \frac{\eta}{|y|}\frac{\mu_R}{P^z},{\mu\over |y|P^z}\Bigr)
 , \nn
\end{align}
其中 $\bar{Z}'_X$ 由傅里叶变换定义：
\begin{align}
\bar{Z}'_X\left( \eta^2,{\mu_R^2\over \mu^2}\right)
\equiv  \int {d\tau\over 2\pi}\ e^{i\eta \tau }\, Z'_X\Bigl(\tau^2,\frac{\mu_R^2}{\mu^2}\Bigr)
  \,.
\end{align}
视方案 $X$ 而定，我们指出稍作修改的 $\bar{Z}'_X$ 定义可能更为合适。

$\overline{\rm MS}$ 方案对重正化\spcorr 的一个不理想之处是它不具有光滑的 $z\to 0$ 极限，因而与``$z=0$ 的定域算符是守恒流''这一事实没有简单联系。为避免这一点，可以简单地使用一个与 $\overline{\rm MS}$ 有简单关系的别的方案，例如把 $C_0(\mu^2z^2)$ 加到 $\overline{\rm MS}$ 重正化常数上。这就去掉了讨厌的 $\ln(\mu^2 z^2)$ 项，得到一个与守恒流光滑衔接的方案。

\begin{figure*}
	\centering
	\includegraphics[width=.49\textwidth]{images/Compare_qPDF_MSbar_Pzycut_TeX.pdf}
	\vspace{-0.3cm}
	\caption{
		$\overline{\rm MS}$ 方案的 PDF $xf_{u-d}$ 与由 $x\, C^{\overline{\rm MS}}(p^z)\otimes f_{u-d}$ 得到的 $\overline{\rm MS}$ 准 PDF，比较了取 $p^z=yP^z$ 与 $p^z=P^z$ 的结果。
	}
	\label{fig:quasi}
\end{figure*}

\begin{figure*}
	\centering
	\includegraphics[width=.49\textwidth]{images/Compare_pPDFvsPDF_TeX.pdf}
	\hspace{0.1cm}
	\includegraphics[width=0.49\textwidth]{images/Compare_pPDFvsPDF_z_TeX.pdf}	\\
	\vspace{-0.3cm}
	\caption{（左）$\overline{\rm MS}$ 方案下 PDF $xf_{u-d}$ 与赝 PDF $x({\cal C}^{\overline{\rm MS}}\otimes f_{u-d})$ 的比较。橙色带与蓝色带表示因子化标度 $\mu=4\,{\rm GeV}$ 变化一倍的结果。（右）同左图，但只显示不同 $z$ 值的中心赝 PDF 曲线。 }
	\label{fig:pseudo}
\end{figure*}

除 $\overline{\rm MS}$ 方案外，准 PDF 还曾被定义为带横向动量截断~\cite{Xiong:2013bka,Lin:2014zya,Ma:2014jla,Alexandrou:2015rja}以及 RI/MOM 方案~\cite{Constantinou:2017sej,Alexandrou:2017huk,Chen:2017mzz,Green:2017xeu,Stewart:2017tvs}。RI/MOM 方案近来引起了强烈兴趣，因为它可以在格点上非微扰实现，所以我们把它作为上述关系的一个显式例子来考察。在该方案中，重正化常数 $Z_{\rm OM}$ 由对一个离壳夸克态中的\spcorr 施加条件来确定：
\begin{align}
&\left.Z^{-1}_{\rm OM}\,
\tilde{Q}(\zeta=zp^z,z^2/a^2,-p^2a^2)\right|_{p^2=-\mu_R^2,p^z=p_R^z}
\nn\\
& \quad = \tilde{Q}^{(0)}_q(zp_R^z,z^2/a^2,z^2\mu_R^2) = e^{-izp_R^z}\,,
\end{align}
其中 $q$ 表示夸克态，$p^\mu$ 是外部动量，上标的``$(0)$''代表树图矩阵元。于是，
\begin{align}
Z_{\rm OM} &= Z_{\rm OM} (zp_R^z, z^2/a^2,a^2\mu_R^2)\,,\nn\\
Z'_{\rm OM}&=Z'_{\rm OM}(zp_R^z, z^2\mu_R^2, \mu_R^2/\mu^2)\,,
\end{align}
并且这里我们定义
\begin{align}
	&\bar{Z}'_{\rm OM}\left(\eta,{\mu_R^2\over (p_R^z)^2},{\mu_R^2\over \mu^2}\right)\nn\\
	&\equiv
   p_R^z \int {dz\over 2\pi}\ e^{i\eta p_R^z z}\ Z'_{\rm OM}(zp_R^z, z^2\mu_R^2, \mu_R^2/\mu^2)\,.
\end{align}
于是式~(\ref{eq:rimomfac})中的匹配系数成为
\begin{align}
& C^{\rm OM}\left(\frac{x}{y}, {\mu_R\over p_R^z},{\mu_R\over \mu}, {\mu\over yP^z}\right)
 \\
=& \int d\eta\ \bar{Z}'_{\rm OM}\left(\eta,{\mu_R^2\over (p_R^z)^2},{\mu_R^2\over\mu^2}\right) C\left({x\over y}-{\eta\over y}{p_R^z\over P^z}, {\mu\over |y|P^z}\right)
 \,. \nn
\end{align}
$\mu_R$ 与 $p_R^z$ 的选择独立于 $\mu$ 与 $P^z$，文献~\cite{Chen:2017mzz,Stewart:2017tvs}采用了 $p_R^z=P^z$。

应当指出，在格点上由于手征对称性的破缺，矢量型夸克 Wilson 线算符 $\tilde{O}_{\gamma^\mu}(z)$ 可以与标量算符 $\tilde{O}_{\bf 1}(z)$ 混合，如文献~\cite{Constantinou:2017sej,Alexandrou:2017huk,Chen:2017mzz,Green:2017xeu,Chen:2017mie}所讨论。考虑混合效应之后，同样的因子化公式仍可应用于来自格点 QCD 的 RI/MOM 准 PDF。
